\section{第~\ref{sec:dendrites}~章　树突数目}

各类的结构和输入数目已通过解剖和电记录测量；估计~\eqref{eq:dendriteload}~是简单的电容器物理；而输入数目的计算解释依赖于理论和模型。

{\footnotesize
\setlength{\tabcolsep}{3pt}\renewcommand{\arraystretch}{1.15}
\begin{longtable}{>{\raggedright}p{0.5cm}>{\raggedright}p{4.0cm}>{\raggedright}p{5.6cm}>{\raggedright}p{2.3cm}>{\raggedright\arraybackslash}p{1.7cm}}
\toprule
编号 & 论断 & 实验及测量内容 & 来源 & 状态 \\
\midrule\endfirsthead
\toprule
编号 & 论断 & 实验及测量内容 & 来源 & 状态 \\
\midrule\endhead
1 & 膜的比电容 $c_m\approx1$~\textmu F/cm$^2$ & 用吸管从神经元胞体上取下膜片，施加电压阶跃，根据电容电流的衰减求出总电容，再除以面积：三类神经元均为 $0.9$~\textmu F/cm$^2$ & \cite{Gentet2000} & 已测量 \\
2 & 充电时间 $t\approx c_mA\Delta V/I$~\eqref{eq:dendriteload}：大神经元需要几十个同时输入，小神经元一个大突触就够 & 按电容器充电定律所作的估计；数字（$20$ 和 $300$~pF，$0.1$ 和几nA）是数量级 & 本章推导 & 理论 \\
3 & 一个巨大突触给出几纳安的电流，立即把小神经元推到阈值 & 在脑片中同时记录末梢和接收者：单个突触的电流为 $2$--$13$~nA，每个输入脉冲都引起阈上响应，$36^\circ$C时延迟约 $0.4$~ms；接收者输出\nt{GLYC} & \cite{Borst1995,BorstSoria2012} & 已测量 \\
4 & 零个树突：触觉和痛觉感觉神经元的脉冲从传感器绕过细胞体进入脊髓 & 结构（在胞体处一分为二的轴突）已由解剖确定。传导不需要胞体的兴奋性，这一点在这类神经元的经过验证的模型中得到证明：没有可兴奋的胞体，脉冲同样可靠地通过分叉 & \cite{AmirDevor2003} & 理论 \\
5 & 一个树突：嗅觉神经元携带一种受体，传递一条输入线路 & 受体基因实验的综述：每个嗅觉神经元只表达大家族中的一个受体基因 & \cite{Mombaerts2004} & 已测量 \\
6 & 两个树突：声音方向检测器中，来自左右耳的输入落在不同的树突上 & 综述汇集的哺乳动物解剖和记录：来自两耳的输入分布在两个相反方向的树突上 & \cite{Grothe2010} & 已测量 \\
7 & 把输入分到两个树突上使“与”更严格 & 在模型中证明：一个树突上的饱和~\eqref{eq:shunt}~使来自一侧的输入无法到达阈值，符合检测得到改善。没有改变输入排布的直接实验 & \cite{AgmonSnir1998} & 理论 \\
8 & 运动学习回路中约 $7\cdot10^{10}$ 个小\nt{GLUT}神经元，每个 $\sim4$ 个输入 & 用各向同性分离法计数细胞：人脑这一部分约有 $69\cdot10^9$ 个神经元。在活体大鼠中记录一个这样的细胞：触碰时来自少数输入的离散电流成簇到来，它们相加时细胞发放 & \cite{Azevedo2009,Chadderton2004} & 已测量 \\
9 & 有 $n$ 个输入的阈值元件最多分开 $P_{\max}\approx2n$ 个随机模式~\eqref{eq:cover}；对运动学习回路的输出神经元，上界 $\sim2\cdot10^5$，现实估计 $\sim4\cdot10^4$ 个对应关系 & 上界是组合几何的经典结果。现实估计是只有正权重并留有抗噪余量的容量理论，应用于这个神经元实测的输入数 & \cite{Cover1965,Brunel2004} & 理论 \\
10 & 输入少的混合器用于扩展输入；“四个”接近最优 & 理论：混合器层给出的表示维数，大约在解剖学上观察到的输入数目处达到最大；扩展的最初假说见经典工作 & \cite{Marr1969,Albus1971,LitwinKumar2017} & 假说 \\
11 & 大树突树：$10^4$--$10^5$ 个输入 & 按电子显微照片计数突触：大鼠运动学习回路的一个输出\nt{GABA}神经元上有 $\approx1.7\cdot10^5$ 个突触；海马\nt{GLUT}神经元上约有 $30\,000$ 个兴奋性输入和 $1\,700$ 个抑制性输入 & \cite{NapperHarvey1988,Megias2001} & 已测量 \\
12 & 大树突树的分支是各自带阈值的独立子电路；神经元是一个小型多层网络 & 在细树突的两个位置点刺激：同一分支上的输入按S形曲线相加，不同分支上的输入线性相加。在人类皮层神经元的树突中发现了分级的钙脉冲，使单个细胞能解决线性不可分的任务。模型：要重现一个神经元需要一个深层网络 & \cite{Polsky2004,Gidon2020,Beniaguev2021} & 已测量 \\
13 & 树突即输出：视网膜\nt{GABA}神经元的每个分支都是独立的方向检测器 & 在单个分支中进行双光子钙成像：分支中的信号对从细胞体向末梢的运动有选择性，而胞体电压没有 & \cite{Euler2002} & 已测量 \\
14 & 输入数目跟随任务（命题~\ref{prop:dendrites}） & 各类的结构已测量（第~3--13行）；输入数目是为了速度或分类而选定的，这一点只有针对个别类别的理论和模型给出 & \cite{LitwinKumar2017,AgmonSnir1998} & 假说 \\
\bottomrule
\end{longtable}}
